%!TEX root = ../main.tex

\newenvironment{abstract}{}{}
\usepackage{abstract}
% \usepackage[T1]{fontenc} % 疑似修改点
% \usepackage{newpxtext,newpxmath} % 疑似修改点
\usepackage{arxiv}
\usepackage[square,numbers]{natbib}






\usepackage[utf8]{inputenc} % allow utf-8 input
\usepackage[T1]{fontenc}    % use 8-bit T1 fonts
% \usepackage{tgpagella}
% \usepackage{mathpazo}
% \usepackage{tgtermes}
% \usepackage{newtxmath}
\usepackage[english]{babel}
\usepackage{dsfont} %
\usepackage{xcolor}

\usepackage{hyperref}
\definecolor{linkgreen}{HTML}{009000}
\definecolor{linkcobalt}{HTML}{3646F0}
\hypersetup{
    colorlinks=true,
    linkcolor=black,
    citecolor=linkgreen,
    filecolor=linkcobalt,
    urlcolor=linkcobalt,
    pdfpagemode=FullScreen,
    }    % hyperlinks
\usepackage{url}            % simple URL typesetting
\usepackage{booktabs}       % professional-quality tables
\usepackage{amsfonts}       % blackboard math symbols
\usepackage{nicefrac}       % compact symbols for 1/2, etc.
\usepackage{microtype}      % microtypography
\usepackage{lipsum}		% Can be removed after putting your text content
\usepackage{graphicx}

\usepackage{doi}
\usepackage{amsmath}
\usepackage{gensymb}
\usepackage{enumitem}
\DeclareMathOperator*{\argmax}{arg\,max}
\DeclareMathOperator*{\argmin}{arg\,min}
\usepackage{pifont}
\usepackage{xparse}
\usepackage{xspace}

\newcommand{\vs}{\textit{vs.}\xspace}
\newcommand{\ie}{\textit{i.e.,}\xspace}
\newcommand{\eg}{\textit{e.g.,}\xspace}
\usepackage{wrapfig}
\usepackage{graphicx}


%\usepackage[landscape,a4paper]{geometry}
\usepackage{colortbl}
\newcommand{\ra}[1]{\renewcommand{\arraystretch}{#1}}
\usepackage{longtable}

\usepackage{titlesec}
% 四级标题
\titleformat{\paragraph}[runin]{\normalsize\bfseries}{\theparagraph}{1em}{}
\titlespacing*{\paragraph}{0pt}{6pt}{4pt}
% 五级标题
\titleformat{\subparagraph}[runin]{\normalsize\itshape}{\thesubparagraph}{1em}{}
\titlespacing*{\subparagraph}{0pt}{4pt}{3pt}
\renewcommand\theparagraph{\thesubsubsection.\arabic{paragraph}}
\renewcommand\thesubparagraph{\theparagraph.\arabic{subparagraph}}
\setcounter{secnumdepth}{4}
\setcounter{secnumdepth}{5} % 允许显示五级标题的序号

%\usepackage[numbers]{natbib}
% \usepackage[margin=1in]{geometry}
\usepackage{tikz}
\usetikzlibrary{trees, tikzmark, shadows, arrows.meta, positioning, calc}
%\usepackage{forest}

% draw the tree
\usepackage[edges]{forest}
\usepackage[normalem]{ulem}
% For timeline
\usepackage[most]{tcolorbox}

% 158 195 231
\definecolor{ai_blue}{RGB}{158, 195, 231}

\definecolor{forestgreen}{RGB}{34,139,34}
\definecolor{steelred}{RGB}{139,30,30}
\newtcolorbox{quotebox}{
    colback=ai_blue!60, 
    colframe=gray!30,
    coltext=black,
    boxrule=0pt,
    arc=4pt,
    left=10pt,
    right=10pt,
    top=10pt,
    bottom=10pt,
    enhanced,
    drop shadow={black!50!white}
}

\usepackage{multicol}
\usepackage{multirow}
\usepackage{cleveref}

\tikzset{%
    parent/.style =          {align=center,text width=2cm,rounded corners=3pt, line width=0.3mm},%, fill=gray!10,draw=gray!80},
    child/.style =           {align=center,text width=2.3cm,rounded corners=3pt, fill=blue!10,draw=blue!80,line width=0.3mm},
    grandchild/.style =      {align=center,text width=2cm,rounded corners=3pt},
    greatgrandchild/.style = {align=center,text width=1.5cm,rounded corners=3pt},
    greatgrandchild2/.style = {align=center,text width=1.5cm,rounded corners=3pt},    
    referenceblock/.style =  {align=center,text width=1.5cm,rounded corners=2pt},
    %%%% Re-define
    % Pretrain Model
    pretrain/.style =           {align=center,text width=2.9cm,rounded corners=3pt, fill=blue!10,draw=blue!80,line width=0.3mm},   
    pretrain_work/.style =           {align=center, text width=2.9cm,rounded corners=3pt, fill=blue!10,draw=blue!0,line width=0.3mm},  
    % Template Mining
    template/.style =           {align=center,text width=2.9cm,rounded corners=3pt, fill=red!10,draw=red!80,line width=0.3mm},   
    template_work/.style =           {align=center,text width=2.9cm,rounded corners=3pt, fill=red!10,draw=red!0,line width=0.3mm},    
    % Answering Mining
    answer/.style =           {align=center,text width=2.9cm,rounded corners=3pt,line width=0.3mm},%, fill= cyan!10,draw= cyan!80,},   
    answer_work/.style =           {align=center,text width=2.9cm,rounded corners=3pt,line width=0.3mm},%, fill= cyan!10,draw= cyan!0,},      
    % Multi-prompt 
    multiple/.style =           {align=center,text width=2.9cm,rounded corners=3pt, fill= orange!10,draw= orange!80,line width=0.3mm},   
    multiple_work/.style =           {align=center,text width=2.9cm,rounded corners=3pt, fill= orange!10,draw= orange!0,line width=0.3mm},        
    % Tuning Strategy
    tuning/.style =           {align=center,text width=2.9cm,rounded corners=3pt, fill= magenta!10,draw= magenta!80,line width=0.3mm},   
    tuning_work/.style =           {align=center,text width=2.9cm,rounded corners=3pt, fill= magenta!10,draw= magenta!0,line width=0.3mm},
    % Assurance
    assurance/.style =           {align=center,text width=2.9cm,rounded corners=3pt, fill= green!10,draw= green!80,line width=0.3mm},   
    assurance/.style =           {align=center,text width=2.9cm,rounded corners=3pt, fill= green!10,draw= green!0,line width=0.3mm},
}


\usepackage{color}			% Color extended names

%DEFINE ENVIRONMENT BLOCK 
% Riddle
\newenvironment{Riddle}{		% Riddle environment
  \begin{center}			% Centering minipage
    \colorbox[HTML]{F8E0E0} {		% Set's the color of minipage
      \begin{minipage}[b]{380px} 	% Starts minipage
 	\textbf{Riddle: }\\ \itshape 	% Set's title and starts italic for text 
      \end{minipage}}			% End minipage
}{\end{center}}				% End Riddle environment


\usepackage{tabularx}
\usepackage{tabularray}
\newcolumntype{C}[1]{>{\centering\arraybackslash}m{#1}}  % centered, auto-wrapping fixed-width column (fraction of \textwidth)
\newcolumntype{L}[1]{>{\raggedright\arraybackslash}m{#1}}  % left-aligned, auto-wrapping fixed-width column
\usepackage{CJKutf8}  % for Chinese text


\usepackage{amssymb}
\usepackage{amsthm}
\usepackage{mathtools}


\tcbuselibrary{theorems,breakable}

% First, adjust the theorem style to match proof spacing
\newtheoremstyle{definition}
  {6pt}        % Space above - matching \topsep6pt from proof
  {6pt}        % Space below - matching \topsep6pt from proof  
  {\normalfont} % Body font
  {}           % Indent amount
  {\bfseries}  % Theorem head font
  {.}          % Punctuation after theorem head
  {.5em}       % Space after theorem head (similar to \labelsep)
  {}           % Theorem head spec

\theoremstyle{definition}
\newtheorem{definition}{Definition}

% Then adjust tcolorbox to preserve these spacings
\definecolor{DefGreen}{HTML}{E8F5E9}
\tcolorboxenvironment{definition}{
  colback=DefGreen, %blue!5!white,
  boxrule=0pt,
  boxsep=0pt,
  % Match internal spacing to typical proof environment
  left=2pt,    % Use same label separation as proof
  right=\labelsep,   % Symmetric padding
  top=12pt,           % Small top padding inside box
  bottom=12pt,        % Small bottom padding inside box
  oversize=2pt,
  sharp corners,
  % Match external spacing to \topsep
  before skip=6pt,   % Same as proof \topsep6pt
  after skip=6pt,    % Same as proof \topsep6pt
  breakable,
}

% Alternative color schemes you might prefer for definitions:
% Option 1: Light gray (more neutral)
% colback=gray!5!white,

% Option 2: Light green (suggesting growth/construction)
% colback=green!5!white,

% Option 3: Light yellow (highlighting important concepts)
% colback=yellow!5!white,




\newcommand{\red}[1]{\textcolor{red}{#1}}
\newcommand{\blue}[1]{\textcolor{red}{#1}}
\newcommand{\bang}[1]{\textcolor{red}{#1}}
\newcommand{\mc}[1]{\textcolor{red}{MC#1}}
\newcommand{\hs}[1]{\textcolor{red}{Huan:#1}}
\newcommand{\kunlun}[1]{{{\textcolor{red}{[Kunlun: #1]}}}}
\newcommand{\xlqi}[1]{{{\textcolor{red}{[Xiaoliang: #1]}}}}
\newcommand{\xf}[1]{{{\textcolor{red}{[Xinfeng: #1]}}}}
\newcommand{\hsr}[1]{{{\textcolor{red}{[Sirui: #1]}}}}
\newcommand{\hbc}[1]{{\textcolor{red}{[haibo: #1]}}}
\newcommand{\hb}[1]{{\textcolor{red}{ #1}}}
\newcommand{\tjhe}[1]{{{\textcolor{red}{[Tanjin: #1]}}}}
\newcommand{\yg}[1]{\textcolor{red}{Yu Gu:#1}}
\newcommand{\jy}[1]{\textcolor{red}{Jiayi:#1}}

\newcommand{\yuyu}[1]{\textcolor{red}{Yuyu:#1}}


\definecolor{customgreen}{RGB}{191,255,191} % BFFFBF in RGB
\definecolor{customblue}{RGB}{228,254,254} % E4FEFE in RGB
\definecolor{customyellow}{RGB}{255,255,191} % FFFFBF in RGB
\definecolor{customred}{RGB}{255,191,191} % FFBFBF in RGB

\definecolor{sealblue}{HTML}{173F67}
\definecolor{sealgold}{HTML}{B08B3E}

% =============
% \DeclareRobustCommand{\FramedGlyph}[1]{%
%   \tikz[baseline=(Letter.base)]{%
%     \node[
%       rounded corners=1.5pt,
%       fill=sealblue,
%       draw=sealgold,
%       double=white,
%       double distance=1.2pt,
%       line width=0.7pt,
%       minimum width=1.02cm,
%       minimum height=1.12cm,
%       inner sep=0pt
%     ] (Frame) {};%    
%     \node[
%       font=\fontsize{29}{29}\selectfont\bfseries,
%       text=white,
%       inner sep=0pt,
%       yshift=-4pt
%     ] (Letter) at (Frame.center) {#1};%
%   }%
% }

% \newcommand{\FramedInitial}[2]{%
%   \lettrine[
%     lines=3,
%     lraise=0.02,
%     lhang=0.00,
%     loversize=0,
%     findent=0.35em,
%     nindent=0.60em,
%     refstring=T
%   ]{\FramedGlyph{#1}}{#2}%
% }

\DeclareRobustCommand{\FramedGlyph}[1]{%
  \raisebox{0.10\baselineskip}{%
    \tikz[baseline=(Letter.base)]{%
      % 固定整个图形的边界，避免 lettrine 误判高度
      \path[use as bounding box] (0,0) rectangle (1.20,1.34);

      % 外框
      \node[
        rounded corners=2pt,
        fill=sealblue,
        draw=sealgold,
        line width=0.9pt,
        minimum width=1.02cm,
        minimum height=1.12cm,
        inner sep=0pt
      ] (Frame) at (0.60,0.67) {};

      % 内部细边框
      \draw[
        white,
        line width=0.45pt
      ]
      ([xshift=2.2pt,yshift=-2.2pt]Frame.north west)
      rectangle
      ([xshift=-2.2pt,yshift=2.2pt]Frame.south east);

      % 字母：向上微调，解决“视觉上偏下”
      \node[
          % font=\initfamily\fontsize{29}{29}\selectfont,
          font=\fontsize{29}{29}\selectfont\bfseries,
        text=white,
        inner sep=0pt,
        yshift=1.6pt
      ] (Letter) at (Frame.center) {#1};
    }%
  }%
}
\newcommand{\FramedInitial}[2]{%
  \lettrine[
    lines=3,
    lraise=0.02,
    lhang=0.00,
    loversize=0,
    findent=0.35em,
    nindent=0.60em,
    refstring=T
  ]{\FramedGlyph{#1}}{#2}%
}